\chapter{引言}\label{s:intro}

世界各地涌现出许多草根团体，
向\grefdex{g:free-range-learner}{散养型学习者}{散养型学习者}
教授编程、网页设计、机器人技术和其他技能，
这些人不在传统课堂里学习。
这些团体存在的意义，是让学习者不必独自摸索；
但讽刺的是，
团体的创办者和老师往往也在自学怎么教书。

其实有更好的办法。
就像了解一些关于病菌和营养的基本常识有助于保持健康一样，
了解一点认知心理学、
教学设计、
包容性
和社区组织方面的知识，
也能让你成为更有效的老师。
本书介绍了一些你现在就能用得上的关键理念，
解释了为什么我们相信它们是对的，
并为你指出其他能帮你走得更远的资源。

\begin{aside}{再利用}
  本书的部分内容最初是为
  \hreffoot{http://carpentries.github.io/instructor-training/}{Software Carpentry 讲师培训项目}
  编写的，
  在 Creative Commons 署名—非商业性使用 4.0 许可协议下
  （\appref{s:license}），
  全书都可以自由分发和再利用。
  你可以在任何课堂（免费或收费）中使用
  \url{http://teachtogether.tech/} 上的在线版本，
  也可以依据\hreffoot{https://zh.wikipedia.org/wiki/合理使用}{合理使用}条款引用其中的简短片段，
  但未经事先许可，不得在商业作品中大量转载。

  我们欢迎各种贡献、更正和建议，
  每出版一个新版本，都会向所有贡献者致谢。
  详情请见\appref{s:joining}，
  行为准则请见\appref{s:conduct}。
\end{aside}

\seclbl{你是谁}{s:intro-audience}

\secref{s:process-personas}介绍了如何弄清你的学习者是谁。
本书面向的四类人都是\grefdex{g:end-user-teacher}{终端用户教师}{终端用户教师}：
教书并不是他们的主业，
他们几乎没有或完全没有教育学背景，
而且可能在体制内课堂之外工作。

\begin{description}

\item[Emily]
  学的是图书情报专业，
  现在在一家小型咨询公司做网页设计和项目管理。
  业余时间，她帮忙为想转行进入科技行业的女性开设网页设计课。
  她正在招募同事，想在本地多开一些课，
  并且想知道如何做出别人也能用的课程，
  以及怎样把一个志愿者教学组织发展起来。

\item[Moshe]
  是一名职业程序员，有两个十几岁的孩子，
  孩子们的学校不开设编程课。
  他自愿每月办一次课后编程俱乐部；
  虽然经常给同事做演示，
  却没有任何课堂教学经验。
  他想学会如何在合理的时间内做出有效的课程，
  也想多了解自主进度的在线课程的利弊。

\item[Samira]
  是机器人专业的本科生，正在考虑毕业后当老师。
  她想在面向同龄人的周末机器人工作坊里帮忙，
  但从未教过课，
  而且\gref{g:impostor-syndrome}{冒充者综合征}很严重。
  她想多了解教育这一行，好判断自己是否适合，
  同时也在寻找能帮她更有效地讲课的具体技巧。

\item[Gene]
  是计算机科学教授。
  他们已经教了六年操作系统本科课程，
  越来越觉得一定有更好的办法。
  他们所在大学的教学发展中心提供的培训，
  只有关于在在线学习管理系统里发布作业和提交成绩的内容，
  所以他们想知道还应该去争取些什么。

\end{description}

这些人有\emph{各种各样的技术背景}
和\emph{一些教学经验}，
但\emph{没有受过教学、课程设计或社区组织方面的正规训练}。
他们大多面对的是\emph{散养型学习者}，
\emph{关注的是青少年和成年人}
而不是儿童；
他们\emph{时间和资源都很有限}。
我们预计这四个人会这样使用本书：

\begin{description}

\item[Emily]
  会和她的志愿者们一起参加每周的在线读书会。

\item[Moshe]
  会在某个周末的一天工作坊里覆盖本书的一部分，
  其余部分自己学。

\item[Samira]
  会在一门一个学期的本科课程里使用本书，配有作业、一个项目和一个期末考试。

\item[Gene]
  会在办公室里或通勤路上自己读这本书，
  一路读，一路希望大学能为高质量教学提供更多支持。

\end{description}

\seclbl{可以读些什么代替}{s:intro-instead}

如果你很忙，或者只想尝尝这本书会讲些什么，
\cite{Brow2018}给出了十条基于证据的计算机教学建议。
你可能还会喜欢：

\begin{itemize}

\item
  \hreffoot{http://carpentries.github.io/instructor-training/}{The Carpentries 讲师培训}，
  本书正是由它衍生而来。

\item
  \cite{Lang2016} 和~\cite{Hust2012}，
  简短易读，
  把你能立刻上手做的事和支撑它们的科研联系了起来。

\item
  \cite{Berg2012,Lemo2014,Majo2015,Broo2016,Rice2018,Wein2018b}
  都充满了可以在课堂上落地的实用建议，
  但等你有了一个框架、能理解这些点子为什么管用之后，它们可能会更有意义。

\item
  \cite{DeBr2015}，
  它通过说明什么不是真的，来解释教育中什么是真的；
  还有~\cite{Dida2016}，
  它把学习理论建立在认知心理学的基础上。

\item
  \cite{Pape1993}，
  它至今仍描绘着一个关于计算机如何能改变教育的鼓舞人心的愿景。
  \hreffoot{https://medium.com/bits-and-behavior/mindstorms-what-did-papert-argue-and-what-does-it-mean-for-learning-and-education-c8324b58aca4}{Amy Ko 的精彩解读}
  对 Papert 思想的总结比我所能做到的更好，
  而~\cite{Craw2010} 则是与前两者相配的一本引人深思的读物。

\item
  \cite{Gree2014,McMi2017,Watt2014} 解释了为什么过去四十年里那么多教育改革都失败了、
  营利性大学是如何利用并加剧社会不平等的，
  以及技术如何一再未能给教育带来革命。

\item
  \cite{Brow2007} 和~\cite{Mann2015}，
  因为你若不去改变我们教书所处的这个系统，就没法教好，
  而单靠你自己又是改变不了的。

\end{itemize}

想要更学术的材料的人，可能也会觉得~\cite{Guzd2015a,Hazz2014,Sent2018,Finc2019,Hpl2018} 颇有收获，
而 \hreffoot{http://computinged.wordpress.com}{Mark Guzdial 的博客}则一直信息量大、发人深思。

\seclbl{致谢}{s:intro-acknowledgments}

没有以下诸位的贡献，就不会有这本书：
Laura Acion、
Jorge Aranda、
Mara Averick、
Erin Becker、
Yanina Bellini Saibene、
Azalee Bostroem、
Hugo Bowne-Anderson、
Neil Brown、
Gerard Capes、
Francis Castro、
Daniel Chen、
Dav Clark、
Warren Code、
Ben Cotton、
Richie Cotton、
Karen Cranston、
Katie Cunningham、
Natasha Danas、
Matt Davis、
Neal Davis、
Mark Degani、
Tim Dennis、
Paul Denny、
Michael Deutsch、
Brian Dillingham、
Grae Drake、
Kathi Fisler、
Denae Ford、
Auriel Fournier、
Bob Freeman、
Nathan Garrett、
Mark Guzdial、
Rayna Harris、
Ahmed Hasan、
Ian Hawke、
Felienne Hermans、
Kate Hertweck、
Toby Hodges、
Roel Hogervorst、
Mike Hoye、
Dan Katz、
Christina Koch、
Shriram Krishnamurthi、
Katrin Leinweber、
Colleen Lewis、
Dave Loyall、
Paweł Marczewski、
Lenny Markus、
Sue McClatchy、
Jessica McKellar、
Ian Milligan、
Julie Moronuki、
Lex Nederbragt、
Aleksandra Nenadic、
Jeramia Ory、
Joel Ostblom、
Elizabeth Patitsas、
Aleksandra Pawlik、
Sorawee Porncharoenwase、
Emily Porta、
Alex Pounds、
Thomas Price、
Danielle Quinn、
Ian Ragsdale、
Erin Robinson、
Rosario Robinson、
Ariel Rokem、
Pat Schloss、
Malvika Sharan、
Florian Shkurti、
Dan Sholler、
Juha Sorva、
Igor Steinmacher、
Tracy Teal、
Tiffany Timbers、
Richard Tomsett、
Preston Tunnell Wilson、
Matt Turk、
Fiona Tweedie、
Martin Ukrop、
Anelda van der Walt、
Stéfan van der Walt、
Allegra Via、
Petr Viktorin、
Belinda Weaver、
Hadley Wickham、
Jason Williams、
Simon Willison、
Karen Word、
John Wrenn、
和 Andromeda Yelton。
我也感谢 Lukas Blakk 设计了标志，
感谢 Shashi Kumar 提供 LaTeX 方面的帮助，
感谢 Markku Rontu 让图表更好看，
还要感谢多年来使用这些材料的所有人。
书中若仍有错误，都由我负责。

\seclbl{练习}{s:intro-exercises}

每一章最后都有各种练习，
标注了建议的形式，
以及面对面进行时通常需要多长时间。
大多数练习也可以换成别的形式来做——尤其是，
如果你在自学这本书，
很多原本为小组设计的练习你依然可以做——
而且你随时可以比建议的时间多花些功夫。

如果你在讲师培训工作坊里使用这些材料，
可以在活动前一两天把下面的练习发给学员，
以便了解他们是谁、怎样最好地帮助他们。
动手之前，请先读一读\secref{s:classroom-prior}里的注意事项。

\exercise{高峰与低谷}{全班}{5}

简要回答以下问题，并与同伴分享。
（如果你按\secref{s:classroom-notetaking}所述在网上一起记笔记，
就把答案写在那里。）

\begin{enumerate}

\item
  你上过的最好的课或工作坊是哪一次？
  它好在哪儿？

\item
  最糟的是哪一次？
  它糟在哪儿？

\end{enumerate}

\exercise{认识你自己}{全班}{10}

与同伴分享以下问题的简要答案。
把答案记下来，方便你在读本书其余部分时回头翻看。

\begin{enumerate}

\item
  你最想教什么？

\item
  你最想教谁？

\item
  你为什么想教书？

\item
  你怎么知道自己教得好？

\item
  关于教与学，你最想了解什么？

\item
  关于教与学，你坚信的一条具体看法是什么？

\end{enumerate}

\exercise{为什么要学编程？}{个人}{20}

政客、商界领袖和教育工作者常说，
人们应该学编程，因为未来的工作会需要它。
然而，
正如 Benjamin Doxtdator \hreffoot{http://www.longviewoneducation.org/field-guide-jobs-dont-exist-yet/}{所指出的}，
这类说法很多都站不住脚。
就算它们是真的，教育也不该是为了让人们胜任未来的工作：
教育应该赋予人们权力，去决定有哪些工作，
并确保这些工作值得去做。
正如\hreffoot{https://computinged.wordpress.com/2017/10/18/why-should-we-teach-programming-hint-its-not-to-learn-problem-solving/}{Mark Guzdial 所指出的}，
学编程的理由其实有很多：

\begin{enumerate}

\item
  理解我们所在的世界。

\item
  研究和理解各种过程。

\item
  能够对影响自己生活的事物提出问题。

\item
  使用一种重要的新读写能力。

\item
  有一种学习艺术、音乐、科学和数学的新方式。

\item
  作为一项职业技能。

\item
  更好地使用计算机。

\item
  作为学习解决问题的一种媒介。

\end{enumerate}

画一个 $3{\times}3$ 的网格，坐标轴标上「低」「中」「高」，
然后把每条理由放进其中一个格子里，
依据是它对你有多重要（X 轴）
以及对你打算教的人有多重要（Y 轴）。

\begin{enumerate}

\item
  哪些点的重要性彼此接近
  （也就是位于网格的对角线上）？

\item
  哪些点彼此不匹配
  （也就是位于非对角线的角上）？

\item
  这应该怎样影响你教什么？

\end{enumerate}
